\section{Цель экспериментальной части}

Экспериментальная часть работы проверяет не только качество итоговых шуток, но и работоспособность всего предложенного конвейера. В частности, эксперименты должны ответить на четыре вопроса:
\begin{enumerate}
    \item можно ли обучать LoRA-адаптеры для MRVF-цели на моделях Qwen3-1.7B и Qwen3-4B;
    \item как MRVF-модели сравниваются с соответствующими базовыми моделями;
    \item влияет ли режим явного рассуждения на качество короткой юмористической генерации;
    \item согласуется ли оптимизация внутренней многореференсной цели с внешней автоматической попарной оценкой.
\end{enumerate}

Эксперименты в данной работе являются предварительными: они предназначены для проверки реализуемости метода и выявления основных проблем постановки. Поэтому результаты следует интерпретировать как диагностические, а не как окончательное утверждение о превосходстве или несостоятельности MRVF для генерации юмора.

\section{Используемые данные}

Во всех экспериментах используется набор данных \texttt{halaction/humor-generation}. Исходная коллекция содержала приблизительно \(664\,000\) текстов шуток. После удаления точных дублей и почти совпадающих текстов текущий корпус содержит приблизительно \(570\,000\) шуток. На его основе были построены векторные представления, группы ключевых слов и многореференсные запросы.

Для обучения MRVF использовалась конфигурация данных \texttt{references}. Каждая строка этого набора содержит запрос, построенный из ключевых слов, и несколько найденных эталонных шуток. Обучение проводилось на части \texttt{train}, а генерация и оценка кандидатов --- на части \texttt{validation}. Разбиение было зафиксировано заранее, чтобы одни и те же запросы могли использоваться для сравнения разных моделей.

Основные объёмные характеристики данных приведены в Таблице~\ref{tab:experiment-data}.

\begin{table}[h]
    \centering
    \begin{tabular}{lr}
        \hline
        Компонент                                & Количество           \\
        \hline
        Исходная коллекция шуток                 & \(\approx 664\,000\) \\
        Корпус после дедупликации                & \(\approx 570\,000\) \\
        Группы ключевых слов                     & \(\approx 1\,500\)   \\
        Многореференсные запросы                 & \(\approx 14\,600\)  \\
        Максимальное число эталонов на запрос    & 5                    \\
        Число эталонов в обучении на один запрос & 2                    \\
        \hline
    \end{tabular}
    \caption{Объёмные характеристики данных, использованных в экспериментах}
    \label{tab:experiment-data}
\end{table}

\section{Обучаемые модели}

В экспериментах использовались две базовые языковые модели семейства Qwen3:
\begin{itemize}
    \item \texttt{Qwen/Qwen3-1.7B};
    \item \texttt{Qwen/Qwen3-4B}.
\end{itemize}

Для обеих моделей обучались LoRA-адаптеры, а не все параметры базовой модели. Такой выбор снижает требования к памяти и делает возможными отладочные запуски на одной видеокарте A100. В обоих случаях обучение выполнялось в формате bfloat16 с включённой экономией памяти через gradient checkpointing.

Основные параметры обучения приведены в Таблице~\ref{tab:mrvf-training-configs}. В таблицу включены только параметры, влияющие на постановку эксперимента; дополнительные инженерные параметры сохранены в конфигурационных файлах проекта.

\begin{table}[h]
    \centering
    \begin{tabular}{lcc}
        \hline
        Параметр                      & Qwen3-1.7B         & Qwen3-4B           \\
        \hline
        Число шагов обучения          & 80                 & 35                 \\
        Скорость обучения             & \(7\cdot 10^{-6}\) & \(5\cdot 10^{-6}\) \\
        Размер группы \(G\)           & 2                  & 2                  \\
        Число эталонов на запрос      & 2                  & 2                  \\
        Макс. длина рассуждения       & 512                & 384                \\
        Макс. длина эталонного ответа & 64                 & 64                 \\
        Целевая функция               & \(\log\)-масса     & \(\log\)-масса     \\
        Нормировка эталонного ответа  & среднее по токенам & среднее по токенам \\
        Нормировка преимущества       & GRPO               & GRPO               \\
        \(\lambda_z\)                 & 1.0                & 1.0                \\
        \(\lambda_{\mathrm{ref}}\)    & 0.1                & 0.1                \\
        LoRA rank                     & 32                 & 16                 \\
        LoRA alpha                    & 64                 & 32                 \\
        Тип чисел                     & bfloat16           & bfloat16           \\
        \hline
    \end{tabular}
    \caption{Основные параметры MRVF-обучения}
    \label{tab:mrvf-training-configs}
\end{table}

Для Qwen3-1.7B использовалась конфигурация \texttt{qwen3-17b-free-think-r5-budgeted.yaml}. Для Qwen3-4B использовалась конфигурация \texttt{qwen3-4b-free-think-r1-budgeted.yaml}. Обе конфигурации используют режим рассуждения Qwen, принудительное закрытие блока рассуждения и MRVF-цель на основе нормированной логарифмической массы.

\section{Сравниваемые варианты моделей}

Для каждой базовой модели сравнивались четыре режима:
\begin{enumerate}
    \item базовая модель без явного рассуждения;
    \item базовая модель с явным рассуждением;
    \item MRVF-модель с явным рассуждением;
    \item MRVF-модель без явного рассуждения.
\end{enumerate}

Такой дизайн нужен, чтобы отделить эффект обучения от эффекта режима рассуждения. Если сравнивать только базовую модель без рассуждения и MRVF-модель с рассуждением, невозможно понять, связано ли изменение качества с обучением или с другим режимом вывода. Поэтому для каждой размерности модели оцениваются все четыре сочетания.

Список сравниваемых вариантов приведён в Таблице~\ref{tab:generation-modes}.

\begin{table}[h]
    \centering
    \begin{tabular}{llll}
        \hline
        Группа & Обозначение                                            & Обучение & Рассуждение \\
        \hline
        1.7B   & \texttt{qwen3-17b-base-nothink-r5eval}                 & нет      & нет         \\
        1.7B   & \texttt{qwen3-17b-base-think-r5eval}                   & нет      & да          \\
        1.7B   & \texttt{qwen3-17b-mrvf-free-think-r5-budgeted}         & MRVF     & да          \\
        1.7B   & \texttt{qwen3-17b-mrvf-free-think-r5-budgeted-nothink} & MRVF     & нет         \\
        4B     & \texttt{qwen3-4b-base-nothink}                         & нет      & нет         \\
        4B     & \texttt{qwen3-4b-base-think}                           & нет      & да          \\
        4B     & \texttt{qwen3-4b-mrvf-r1-budgeted-think}               & MRVF     & да          \\
        4B     & \texttt{qwen3-4b-mrvf-r1-budgeted-nothink}             & MRVF     & нет         \\
        \hline
    \end{tabular}
    \caption{Сравниваемые варианты моделей}
    \label{tab:generation-modes}
\end{table}

\section{Протокол генерации кандидатов}

Кандидатные ответы генерировались на валидационной части набора \texttt{references}. Для каждого запроса модель должна была вернуть одну короткую шутку по заданным ключевым словам. В режимах без рассуждения генерация ограничивалась коротким ответом. В режимах с рассуждением модель сначала получала отдельный бюджет на блок рассуждения, после чего генерировала короткий финальный ответ.

Для Qwen3-1.7B использовалось 200 валидационных запросов. После фильтрации непустых ответов, общих для всех четырёх вариантов, осталось 184 запроса. Для Qwen3-4B использовалось 100 валидационных запросов. Параметры генерации приведены в Таблице~\ref{tab:generation-protocol}.

\begin{table}[h]
    \centering
    \begin{tabular}{lccc}
        \hline
        Группа & Режим           & Бюджет рассуждения & Бюджет финального ответа \\
        \hline
        1.7B   & без рассуждения & --                 & 48 токенов               \\
        1.7B   & с рассуждением  & 2048 токенов       & 48 токенов               \\
        4B     & без рассуждения & --                 & 48 токенов               \\
        4B     & с рассуждением  & 1536 токенов       & 48 токенов               \\
        \hline
    \end{tabular}
    \caption{Параметры генерации кандидатных ответов}
    \label{tab:generation-protocol}
\end{table}

Во всех режимах с рассуждением использовалось принудительное закрытие блока рассуждения перед генерацией финального ответа. Это было необходимо для того, чтобы отделить внутреннее рассуждение модели от текста, который затем сравнивается моделью-оценщиком.

\section{Протокол автоматической оценки}

Автоматическая оценка проводилась попарно. Для каждого запроса и каждой пары моделей формировалось сравнение двух ответов. Порядок левого и правого ответа случайно перемешивался с фиксированным зерном случайности, чтобы уменьшить позиционное смещение. Модель-оценщик получала исходный запрос и два кандидатных ответа, после чего выбирала более удачную шутку.

Для четырёх моделей число пар моделей равно
\[
    {4 \choose 2}=6.
\]
Поэтому для \(N\) общих запросов число попарных сравнений равно
\[
    6N.
\]
В эксперименте с Qwen3-1.7B использовалось \(N=184\) общих непустых запросов, что даёт \(1104\) сравнения. В эксперименте с Qwen3-4B использовалось \(N=100\) запросов, что даёт \(600\) сравнений. Эти объёмы приведены в Таблице~\ref{tab:evaluation-volume}.

\begin{table}[h]
    \centering
    \begin{tabular}{lrrr}
        \hline
        Группа     & Число моделей & Число запросов & Число сравнений \\
        \hline
        Qwen3-1.7B & 4             & 184            & 1104            \\
        Qwen3-4B   & 4             & 100            & 600             \\
        \hline
    \end{tabular}
    \caption{Объём автоматической попарной оценки}
    \label{tab:evaluation-volume}
\end{table}

По результатам попарных сравнений для каждой модели вычислялись число побед, число поражений, доля побед и оценка Брэдли--Терри. Эти величины используются в следующей главе для представления итоговых результатов.

\section{Сохранение артефактов и воспроизводимость}

Для воспроизводимости экспериментов сохранялись несколько типов артефактов:
\begin{itemize}
    \item конфигурации обучения;
    \item логи обучения и выборки промежуточных ответов;
    \item LoRA-адаптеры обученных моделей;
    \item сгенерированные кандидатные ответы;
    \item результаты попарной оценки;
    \item агрегированные таблицы лидеров.
\end{itemize}

LoRA-адаптеры были опубликованы отдельно от базовых моделей в репозиториях \texttt{halaction/Qwen3-1.7B-humor-lora} и \texttt{halaction/Qwen3-4B-humor-lora}. Кандидатные ответы, результаты оценки и таблицы лидеров были добавлены в набор данных \texttt{halaction/humor-generation}. Такая структура позволяет отдельно воспроизводить обучение, генерацию, оценку и агрегацию результатов.

\section{Статус текущего экспериментального прогона}

На момент написания этой версии текста был выполнен полный тестовый прогон: обучение MRVF-адаптеров, генерация кандидатов, загрузка контрольных точек, попарная оценка и построение таблиц лидеров. Однако этот прогон не рассматривается как окончательная оценка качества моделей.

После анализа сгенерированных ответов была обнаружена проблема в шаблоне генерации: часть ответов содержала пояснения, повторяющиеся фрагменты или не отделяла финальную шутку от рассуждения. Для задачи короткой генерации юмора это существенно искажает сравнение, поскольку модель-оценщик сравнивает не только юмористическую идею, но и формат ответа. Поэтому текущий прогон используется как проверка работоспособности конвейера, а не как финальный результат по качеству.

После исправления шаблона генерации кандидаты и попарная оценка должны быть пересчитаны. В следующей главе таблицы результатов представлены в структуре, рассчитанной на замену предварительных значений окончательными после повторного прогона.
